\question\TFQuestion{T}{If populations of a single species are geographically isolated for a \textit{long} period of time, then they might evolve biological reproductive isolation.}

\question\TFQuestion{T}{The species is the \emph{largest} unit within which gene flow can readily occur.}

\question\TFQuestion{F}{Homologies evolve by natural selection. Analogies do not.}

\question\TFQuestion{F}{The basal group is the common ancestor of all other organisms on the phylogenetic tree.}

\question\TFQuestion{F}{At least one prezygotic reproductive barrier must evolve for speciation to occur.}

\question\TFQuestion{F}{A monophyletic group includes the ancestor and most or all of its descendants.}

\question\TFQuestion{F}{Reinforcement occurs between allopatric species.}

\question\TFQuestion{T}{The best classification system is that which most closely reflects evolutionary history.}

\question\TFQuestion{F}{Analogies can be used to test hypotheses if two organisms do \emph{not} share a common ancestor.}

\question\TFQuestion{T}{Both intrasexual and intersexual selection can cause the evolution of sexual dimorphism.}

\question\TFQuestion{T}{Intrasexual selection acts directly on relative fitness.}

